% GENERATED FILE. Edit canonical lessons, then run npm run generate:books.
% canonical-source-hash: fnv1a64:8f63e16eb8ce7c81
% canonical-lessons: HI-C507-sahkarmi, HI-C507-afsar, HI-C507-ustad, HI-C507-kausal, HI-C507-pratibha

\chapter{People at Work}
\label{ch:hi-people-at-work}
\hlchaptermodality{507}

\begin{chapteropening}
\textbf{By the end of this chapter:} \emph{I can talk about the people I work with and what they are good at.}

Complete the last lesson of chapter 507: I can talk about the people I work with and what they are good at.
\end{chapteropening}

\section[\texorpdfstring{\hi{सहकर्मी} (sahkarmī)}{सहकर्मी (sahkarmī)}]{\hi{सहकर्मी} (sahkarmī) — a colleague; a co-worker}

\label{lesson:HI-C507-sahkarmi}



\begin{quote}
Before the new one: say the Hindi for to shape, then the Hindi for a mechanic.
\end{quote}

\subsection*{You'll want to know: \hi{सहकर्मी}}
\textbf{\hi{सहकर्मी}} — \emph{sahkarmī} — \textquotedblleft{}a colleague; a co-worker\textquotedblright{}.

\begin{quote}
\hi{मीरा} \hi{अपनी} \hi{सहकर्मी} \hi{के} \hi{साथ} \hi{दोपहर} \hi{का} \hi{खाना} \hi{खाती} \hi{है।}
\end{quote}

\begin{quote}
\emph{Mīrā apnī sahkarmī ke sāth dopahar kā khānā khātī hai.}
\end{quote}

\begin{quote}
Mira has lunch with her colleague.
\end{quote}

Masculine or feminine, by the person meant.

\begin{cousinweb}
\textbf{\hi{सहकर्मी}} (\emph{sahkarmī}) is a modern compound of two Sanskrit words, \emph{saha}, \textquotedblleft{}with\textquotedblright{}, and \emph{karmin}, \textquotedblleft{}active, busy\textquotedblright{} (Monier-Williams): one who works alongside you.
\end{cousinweb}

\begin{grammarlens}[title={What you've built}]
One more word for describing what you lived and felt.
\end{grammarlens}

\subsection*{Your turn}
\begin{itemize}
\raggedright
  \item \emph{Say it:} \emph{sahkarmī}
  \item \emph{Say it:} \emph{sahkarmī}, once more
  \item \emph{Say it:} \emph{Mīrā apnī sahkarmī ke sāth dopahar kā khānā khātī hai.} — Mira has lunch with her colleague.
  \item \emph{Recall:} say the Hindi for to shape, then the Hindi for a mechanic, then say \emph{sahkarmī} again
\end{itemize}

\begin{tcolorbox}[breakable,colback=teal!4,colframe=teal!35!black,title={Before you move on}]
What does \hi{सहकर्मी} mean? (\textquotedblleft{}A colleague; a co-worker\textquotedblright{}.) Say it once more.
\end{tcolorbox}
\section[\texorpdfstring{\hi{अफ़सर} (afsar)}{अफ़सर (afsar)}]{\hi{अफ़सर} (afsar) — an officer; an official}

\label{lesson:HI-C507-afsar}



\begin{quote}
Before the new one: say the Hindi for a mechanic, then the Hindi for a colleague.
\end{quote}

\subsection*{You'll want to know: \hi{अफ़सर}}
\textbf{\hi{अफ़सर}} — \emph{afsar} — \textquotedblleft{}an officer; an official\textquotedblright{}.

\begin{quote}
\hi{अफ़सर} \hi{ने} \hi{हर} \hi{मज़दूर} \hi{से} \hi{एक} \hi{सवाल} \hi{पूछा।}
\end{quote}

\begin{quote}
\emph{Afsar ne har mazdūr se ek savāl pūchā.}
\end{quote}

\begin{quote}
The officer asked every labourer a question.
\end{quote}

Masculine or feminine, by the person meant.

\begin{cousinweb}
\textbf{\hi{अफ़सर}} (\emph{afsar}) is usually traced to English \emph{officer}, reshaped to fit Hindi sounds.
\end{cousinweb}

\begin{grammarlens}[title={What you've built}]
One more word for describing what you lived and felt.
\end{grammarlens}

\subsection*{Your turn}
\begin{itemize}
\raggedright
  \item \emph{Say it:} \emph{afsar}
  \item \emph{Say it:} \emph{afsar}, once more
  \item \emph{Say it:} \emph{Afsar ne har mazdūr se ek savāl pūchā.} — The officer asked every labourer a question.
  \item \emph{Recall:} say the Hindi for a mechanic, then the Hindi for a colleague, then say \emph{afsar} again
\end{itemize}

\begin{tcolorbox}[breakable,colback=teal!4,colframe=teal!35!black,title={Before you move on}]
What does \hi{अफ़सर} mean? (\textquotedblleft{}An officer; an official\textquotedblright{}.) Say it once more.
\end{tcolorbox}
\section[\texorpdfstring{\hi{उस्ताद} (ustād)}{उस्ताद (ustād)}]{\hi{उस्ताद} (ustād) — a master (of a craft or an art); a teacher}

\label{lesson:HI-C507-ustad}



\begin{quote}
Before the new one: say the Hindi for a colleague, then the Hindi for an officer.
\end{quote}

\subsection*{You'll want to know: \hi{उस्ताद}}
\textbf{\hi{उस्ताद}} — \emph{ustād} — \textquotedblleft{}a master (of a craft or an art); a teacher\textquotedblright{}.

\begin{quote}
\hi{मेरे} \hi{उस्ताद} \hi{ने} \hi{मुझे} \hi{लकड़ी} \hi{का} \hi{काम} \hi{सिखाया।}
\end{quote}

\begin{quote}
\emph{Mere ustād ne mujhe lakṛī kā kām sikhāyā.}
\end{quote}

\begin{quote}
My master taught me woodwork.
\end{quote}

Masculine.

\begin{cousinweb}
\textbf{\hi{उस्ताद}} (\emph{ustād}) is usually traced to Persian \emph{ustād}, \textquotedblleft{}a master\textquotedblright{}. In Hindi it is the master of a craft or an art, and the teacher who passes it on.
\end{cousinweb}

\begin{grammarlens}[title={What you've built}]
One more word for describing what you lived and felt.
\end{grammarlens}

\subsection*{Your turn}
\begin{itemize}
\raggedright
  \item \emph{Say it:} \emph{ustād}
  \item \emph{Say it:} \emph{ustād}, once more
  \item \emph{Say it:} \emph{Mere ustād ne mujhe lakṛī kā kām sikhāyā.} — My master taught me woodwork.
  \item \emph{Recall:} say the Hindi for a colleague, then the Hindi for an officer, then say \emph{ustād} again
\end{itemize}

\begin{tcolorbox}[breakable,colback=teal!4,colframe=teal!35!black,title={Before you move on}]
What does \hi{उस्ताद} mean? (\textquotedblleft{}A master (of a craft or an art); a teacher\textquotedblright{}.) Say it once more.
\end{tcolorbox}
\section[\texorpdfstring{\hi{कौशल} (kauśal)}{कौशल (kauśal)}]{\hi{कौशल} (kauśal) — skill; know-how (formal)}

\label{lesson:HI-C507-kausal}



\begin{quote}
Before the new one: say the Hindi for an officer, then the Hindi for a master (of a craft or an art).
\end{quote}

\subsection*{You'll want to know: \hi{कौशल}}
\textbf{\hi{कौशल}} — \emph{kauśal} — \textquotedblleft{}skill; know-how (formal)\textquotedblright{}.

\begin{quote}
\hi{हर} \hi{कारीगर} \hi{के} \hi{पास} \hi{अपना} \hi{कौशल} \hi{होता} \hi{है।}
\end{quote}

\begin{quote}
\emph{Har kārīgar ke pās apnā kauśal hotā hai.}
\end{quote}

\begin{quote}
Every craftsman has his own skill.
\end{quote}

Masculine.

\begin{cousinweb}
\textbf{\hi{कौशल}} (\emph{kauśal}) is a tatsama, taken whole from Sanskrit \emph{kauśala}, \textquotedblleft{}well-being; skilfulness, cleverness\textquotedblright{} (Monier-Williams).
\end{cousinweb}

\begin{grammarlens}[title={What you've built}]
One more word for describing what you lived and felt.
\end{grammarlens}

\subsection*{Your turn}
\begin{itemize}
\raggedright
  \item \emph{Say it:} \emph{kauśal}
  \item \emph{Say it:} \emph{kauśal}, once more
  \item \emph{Say it:} \emph{Har kārīgar ke pās apnā kauśal hotā hai.} — Every craftsman has his own skill.
  \item \emph{Recall:} say the Hindi for an officer, then the Hindi for a master (of a craft or an art), then say \emph{kauśal} again
\end{itemize}

\begin{tcolorbox}[breakable,colback=teal!4,colframe=teal!35!black,title={Before you move on}]
What does \hi{कौशल} mean? (\textquotedblleft{}Skill; know-how (formal)\textquotedblright{}.) Say it once more.
\end{tcolorbox}
\section[\texorpdfstring{\hi{प्रतिभा} (pratibhā)}{प्रतिभा (pratibhā)}]{\hi{प्रतिभा} (pratibhā) — talent; a gift}

\label{lesson:HI-C507-pratibha}



\begin{quote}
Before the new one: say the Hindi for a master (of a craft or an art), then the Hindi for skill.
\end{quote}

\subsection*{You'll want to know: \hi{प्रतिभा}}
\textbf{\hi{प्रतिभा}} — \emph{pratibhā} — \textquotedblleft{}talent; a gift\textquotedblright{}.

\begin{quote}
\hi{मीरा} \hi{में} \hi{गाने} \hi{की} \hi{बड़ी} \hi{प्रतिभा} \hi{है।}
\end{quote}

\begin{quote}
\emph{Mīrā meṁ gāne kī baṛī pratibhā hai.}
\end{quote}

\begin{quote}
Mira has a great gift for singing.
\end{quote}

Feminine.

\begin{cousinweb}
\textbf{\hi{प्रतिभा}} (\emph{pratibhā}) is a tatsama, taken whole from Sanskrit \emph{pratibhā}, \textquotedblleft{}light; intelligence, genius\textquotedblright{} (Monier-Williams).
\end{cousinweb}

\begin{grammarlens}[title={What you've built}]
One more word for describing what you lived and felt.
\end{grammarlens}

\subsection*{Your turn}
\begin{itemize}
\raggedright
  \item \emph{Say it:} \emph{pratibhā}
  \item \emph{Say it:} \emph{pratibhā}, once more
  \item \emph{Say it:} \emph{Mīrā meṁ gāne kī baṛī pratibhā hai.} — Mira has a great gift for singing.
  \item \emph{Recall:} say the Hindi for a master (of a craft or an art), then the Hindi for skill, then say \emph{pratibhā} again
\end{itemize}

\begin{tcolorbox}[breakable,colback=teal!4,colframe=teal!35!black,title={Before you move on}]
What does \hi{प्रतिभा} mean? (\textquotedblleft{}Talent; a gift\textquotedblright{}.) Say it once more.
\end{tcolorbox}
